% !TeX root = ../main.tex
% ============================================================
%  tables/tab_instances.tex
% ============================================================

\renewcommand{\arraystretch}{1.2}
\begin{tabular}{@{}llp{0.52\linewidth}@{}}
  \toprule
  \textbf{Code} & \textbf{Values} & \textbf{Implemented meaning} \\
  \midrule
  $I$ & 17, 25
    & Total entity nodes in
      $\setS\cup\setF\cup\setW\cup\setR\cup\mathrm{AIR}\cup\mathrm{PORT}$.
      I17 uses 3 suppliers, 3 factories, 3 warehouses, 4 retailers,
      2 airports, and 2 seaports; I25 increases warehouses to 9 and
      airports to 4. \\
  $M$ & 4, 8
    & Material/product scale. M4 uses 2 raw materials and 1 final
      product; M8 uses 4 raw materials and 2 final products. \\
  $T$ & 3, 6, 12
    & Number of planning periods in the generated instance. \\
  Modes & \makecell[l]{mrd1--mrd2 ($\setMrd$),\\ma1--ma2 ($\setMa$),\\ms1--ms2 ($\setMs$)}
    & Common transport catalogue with two alternatives per family,
      following Tautenhain et al. The road modes ($\setMrd$) are
      company-owned truck alternatives; the air ($\setMa$) and sea
      ($\setMs$) modes are outsourced carrier alternatives. \\
  $\taxindex$ & 0,\ldots,20
    & Discrete tax index, used only in the leader's formulation and the
      \DBP (\Cref{sec:model:variables}). The economically meaningful
      rate is $\tau=10\,\taxindex$ USD/t\,CO\textsubscript{2}e, so the
      implemented domain is $[0,200]$ USD/t (justified in
      \Cref{app:data:tax_range}); all reported results use $\tau$. \\
  \bottomrule
\end{tabular}
